\begin{resumo}
O presente documento estabelece a especificação formal do Projeto de Banco de Dados e Modelagem Entidade-Relacionamento do Sistema Inteligente de Gestão de Eventos Adversos baseado na metodologia \textit{Global Trigger Tool} (SIGEA-GTT), desenvolvido no âmbito da Universidade Federal de Sergipe (UFS). O sistema opera sobre a tríplice fronteira de dados composta por controle estrito de identidades institucionais, taxonomia clínica padronizada de eventos adversos (IHI e NCC MERP) e gestão pedagógica de turmas, avaliações e prontuários simulados. A arquitetura de dados adota o modelo híbrido relacional-documental implementado no SGBD PostgreSQL 16 LTS, empregando tabelas estritamente normalizadas (3FN/BCNF) com chaves universais UUID v4 para integridade relacional, combinadas com estruturas semiestruturadas \texttt{JSONB} para representação de prontuários médicos de alta fidelidade e artefatos de ferramentas da qualidade (Ishikawa, 5W2H, GUT e PDCA). Apresentam-se o Modelo Entidade-Relacionamento Conceitual fundamentado na notação de Peter Chen, o Diagrama Lógico Relacional na notação Pé de Galinha (\textit{Crow's Foot}), o dicionário de dados exaustivo de todas as oito entidades físicas, os esquemas formais \textit{JSON Schema}, as políticas de indexação B-Tree e GIN, as regras de integridade de domínio e a matriz de rastreabilidade integral frente aos requisitos do sistema, garantindo isolamento ético e conformidade com a Lei Geral de Proteção de Dados (LGPD).

\vspace{\onelineskip}
\noindent
\textbf{Palavras-chave}: Banco de Dados Relacional. Modelo Entidade-Relacionamento. PostgreSQL. JSONB. Global Trigger Tool. Engenharia de Software.
\end{resumo}
